%%
%% Template file for conference poster
%% TSD2011
%% Technologies for the Sustainable Development
%% 2011
%% 
%% Based on the examples of the beamerposter package from the
%% I6 Informatik Institute at RWTH-Aachen University.

\documentclass[final,hyperref={pdfpagelabels=false}]{beamer}
\mode<presentation>{
  \usetheme{tsd2011}
}

\usepackage{grffile}

%%%%%%%%%%%%%%%%%%%%%%%%%%%%%%%%%%%%%%%%%%%%%%%%%%%%%%%%%%%%%%%%%%%%%%%%%%%%%%%
%%     Please complete the following commands with the corresponding data    %%
%%%%%%%%%%%%%%%%%%%%%%%%%%%%%%%%%%%%%%%%%%%%%%%%%%%%%%%%%%%%%%%%%%%%%%%%%%%%%%%

\title{\huge $<$Poster Title Here$>$}
\author{$<$author or authors$>$}
\institute[$<$Short institute$>$]{$<$Department or Lab, University, City$>$}
\date[Feb. 3-4, 2011]{Feb. 3-4, 2011}
 
\tsdaddress{\insertshortinstitute}
\tsdemail{$<$your e-mail$>$}
\tsdurl{$<$your URL (like http://tsd2011.ietec.org)$>$}
\tsdlogofile{logo} %% the extension is automatically determined!

\usepackage[english]{babel}
\usepackage[utf8]{inputenc}
\usepackage{amsmath,amsthm, amssymb, latexsym}
\usepackage[orientation=portrait,size=a0,scale=1.4,debug]{beamerposter}

\usepackage{array,booktabs,tabularx}
\newcolumntype{Z}{>{\centering\arraybackslash}X} % centered tabularx columns

\newcommand{\pphantom}{\textcolor{red}} % phantom introduces a vertical space
                                        % in p formatted table columns??!!

\listfiles
\graphicspath{{fig/}}
 

%%%%%%%%%%%%%%%%%%%%%%%%%%%%%%%%%%%%%%%%%%%%%%%%%%%%%%%%%%%%%%%%%%%%%%%%%%%%%%%
\newlength{\columnheight}
\setlength{\columnheight}{105cm}


%%%%%%%%%%%%%%%%%%%%%%%%%%%%%%%%%%%%%%%%%%%%%%%%%%%%%%%%%%%%%%%%%%%%%%%%%%%%%%%
\begin{document}
\begin{frame}
  \begin{columns}
    % ---------------------------------------------------------%
    % Set up first column 
    \begin{column}{.49\textwidth}
      \begin{beamercolorbox}[center,wd=\textwidth]{postercolumn}
        % tweaks the width, makes a new \textwidth for the column
        \begin{minipage}[T]{.95\textwidth}  
          % must be some better way to set the the height, width and textwidth
          % simultaneously.
          % Since all columns are the same length, it is all nice and tidy.
          % You have to get the height empirically
          \parbox[t][\columnheight]{\textwidth}{
            % ---------------------------------------------------------%
            % fill each column with content
            {\large{\textbf{General remarks}}}
            \begin{itemize}
            \item The size of the poster must be A0, i.e. 841\,mm$\times$1189\,mm
            \item The poster should make emphasis on your results, and should expose this
              results in a visual way through images, figures, charts, plots, equations, tables,
              etc.
            \item The LaTeX style takes care of most formatting issues, but if you are using 
              other templates, then use sans-serif fonts (like arial or helvetica) with the
              following sizes
              \begin{itemize}
              \item Title: 80\,pt bold
              \item Authors: 44\,pt regular
              \item Institute: 44\,pt regular
              \item Block title: 40\,pt bold
              \item Block text: 36\,pt regular
              \end{itemize}
            \end{itemize}

            \vfill

            \begin{block}{Abstract}
              The abstract could be the same as in the final camera-ready paper, but since
              it is recommended to reduce the amount of text in a poster, you could summarize even
              more your abstract in the poster.
            \end{block}

            \vfill

            \begin{block}{Introduction}
              \begin{itemize}
              \item Present your context
                \begin{itemize}
                \item rely on a solid bibliographical research
                \end{itemize}
              \item Synthesize your problem
                \begin{itemize}
                \item maybe your objectives here
                \item but all the previious is just a suggestion
                \end{itemize}
              \end{itemize}              
            \end{block}
            \vfill
            \begin{block}{Global warming}
              \begin{columns}
                \begin{column}{.65\textwidth}
                  \begin{itemize}
                  \item A second box with a more detailed problem description
                    is suggested
                    \begin{itemize}
                    \item Scenario A
                    \item Scenerio B
                    \end{itemize}
                  \item but they do not consider
                    \begin{itemize}
                    \item Condition A, or
                    \item \alert{Strong} condition B
                    \end{itemize}
                  \end{itemize}
                \end{column}
                \begin{column}{.34\textwidth}
                  \centering
                  \framebox[0.9\textwidth][c]{Some nice figure}
                \end{column}
              \end{columns}
              \vskip-1ex
            \end{block}
            \vfill
            \begin{block}{Maybe some theoretical background}
              \begin{itemize}
              \item Remember that the poster should be clear
                \begin{itemize}
                \item so do not use too much text
                \item but insert as many images you can
                \end{itemize}
              \item You may want to insert a self explanatory equation like:
                \begin{align*}
                  f_{pa}(A_\mathbf{u},\mathcal{G}) &=
                  \sum_{\langle\underline{\mathcal{I}}_i,
                               \underline{\mathcal{S}}_i\rangle\in\mathcal{G}}
                  \frac{\displaystyle\sum_{\underline{\mathcal{R}}_j\in\underline{\mathcal{S}}_i}
                    \left|
                      \mathcal{P}(\underline{\mathcal{R}}_{j},
                                  A_\mathbf{u}(\underline{\mathcal{I}}_i))
                    \right|}
                  {\displaystyle
                    \sum_{\underline{\mathcal{R}}_{j} \in \underline{\mathcal{S}}_i}|
                    \underline{\mathcal{R}}_{j}|}\\
                  \mathcal{P}(\underline{\mathcal{R}}_{j},\mathcal{S}_i) &=
                  \Big\{         
                  \mathbf{p} \mid 
                  \mathbf{p}\!\in\! \underline{\mathcal{R}}_{j}
                  \wedge 
                  \mathbf{p}\!\in\!\mathcal{R}_{k}, 
                  \mathcal{R}_{k}\!\in\!\mathcal{S}_i,
                  j=\arg\max_l | \mathcal{R}_{k}{\cap}\underline{\mathcal{R}}_{l} |
                  \Big\}
                \end{align*}
                (just kidding... some explanation should be here)
              \item Or maybe another figure
              \end{itemize}
            \end{block}
            \vfill
            \begin{block}{Experimental setup}
              \begin{itemize}
              \item The experimental setup used to produce your results should
                be well described
              \item Images with you objects could be used
                \begin{center}
                  \includegraphics[width=0.9\textwidth]{\thetsdlogofile}
                \end{center}
              \end{itemize}
            \end{block}
          } % end of parbox used to force a the column height
        \end{minipage}
      \end{beamercolorbox}
    \end{column}
    % ---------------------------------------------------------%
    % end the column

    % ---------------------------------------------------------%
    % Set up a column 
    \begin{column}{.49\textwidth}
      \begin{beamercolorbox}[center,wd=\textwidth]{postercolumn}
        \begin{minipage}[T]{.95\textwidth} % as before
          \parbox[t][\columnheight]{\textwidth}{ % as before
            % ---------------------------------------------------------%
            % fill each column with content
            
            \begin{block}{Databases}
              \begin{columns}
                \begin{column}{.6\textwidth}
                  \begin{itemize}
                  \item Maybe you used different databases for your experiments
                    \begin{itemize}
                    \item for one set of conditions
                    \item or for another setup
                    \end{itemize}
                  \item Some samples could be presented on the right
                  \end{itemize}
                \end{column}
                \begin{column}{.39\textwidth}
                  \includegraphics[width=0.22\linewidth]{logo}
                  \-
                  \includegraphics[width=0.22\linewidth]{logo}
                  \-
                  \includegraphics[width=0.22\linewidth]{logo}
                  \-
                  \includegraphics[width=0.22\linewidth]{logo}
                  \vskip3ex

                  \includegraphics[width=0.22\linewidth]{logo}
                  \-
                  \includegraphics[width=0.22\linewidth]{logo}
                  \-
                  \includegraphics[width=0.22\linewidth]{logo}
                  \-
                  \includegraphics[width=0.22\linewidth]{logo}
                \end{column}
              \end{columns}
            \end{block}
            \vfill
            \begin{block}{Results 1}
              \begin{itemize}
              \item Results is the most important part of the poster
              \item You will use tables and figures to present them
              \end{itemize}
              \vskip-0.5ex
              \begin{table}
                \centering
                \small
                \begin{tabular}{clcrrr}
                  \toprule 
                  Group &  Category  & Resolution & \multicolumn{3}{c @{}}{Error Rates [\%]} \\
                  \cmidrule(l){4-6}
                  &  &  & \multicolumn{1}{c}{Min} & \multicolumn{1}{c}{Max} & \multicolumn{1}{c @{}}{Avrg} \\
                  \midrule
                  A  & Blue   & 160 $\times$ 120 & 40 &  60 & 50  \\
                  B  & Yellow & 160 $\times$ 120 & 45 &  55 & 50  \\
                  \addlinespace
                  C  & Red    & 160 $\times$ 120 & 10 &  30  & 20  \\  
                  D  & Red    & 320 $\times$ 240 & 10 &  31  & 18  \\
                  E  & Red    & 640 $\times$ 320 & 7  &  28  & 17  \\
                  \addlinespace
                  F  & Green  & 160 $\times$ 120 & 5  &  10  & 7   \\ 
                  G  & Green  & 320 $\times$ 240 & 4  &  11  & 6   \\ 
                  H  & Green  & 640 $\times$ 320 & 3  &  10  & 5   \\
                  \bottomrule
                \end{tabular}
              \end{table}

              \vskip1ex
              \begin{itemize}
              \item Maybe another table here
              \end{itemize}
            \end{block}
            \vfill
            \begin{block}{Results 2}
              \begin{columns}
                \begin{column}{.5\textwidth}
                  \begin{itemize}
                  \item Maybe here a comparison with other methods in the literature
                  \item A table here could be nice
                    \vskip-0.5ex
                    \begin{table}
                      \centering
                      \small
                      \begin{tabular}{clcrrr}
                        \toprule 
                        Group &  Category  & Resolution & \multicolumn{3}{c @{}}{Error Rates [\%]} \\
                        \cmidrule(l){4-6}
                        &  &  & \multicolumn{1}{c}{Min} & \multicolumn{1}{c}{Max} & \multicolumn{1}{c @{}}{Avrg} \\
                        \midrule
                        A  & Blue   & 160 $\times$ 120 & 40 &  60 & 50  \\
                        B  & Yellow & 160 $\times$ 120 & 45 &  55 & 50  \\
                        \addlinespace
                        C  & Red    & 160 $\times$ 120 & 10 &  30  & 20  \\  
                        D  & Red    & 320 $\times$ 240 & 10 &  31  & 18  \\
                        E  & Red    & 640 $\times$ 320 & 7  &  28  & 17  \\
                        \addlinespace
                        F  & Green  & 160 $\times$ 120 & 5  &  10  & 7   \\ 
                        G  & Green  & 320 $\times$ 240 & 4  &  11  & 6   \\ 
                        H  & Green  & 640 $\times$ 320 & 3  &  10  & 5   \\
                        \bottomrule
                      \end{tabular}
                    \end{table}

                    \vskip1ex
                    
                  \end{itemize}
                \end{column}
                \begin{column}{.49\textwidth}                
                  \vskip-3ex
                  \begin{itemize}
                  \item One example figure of their method

                      \includegraphics[width=0.9\linewidth]{logo}

                  \item Another example figure with our method

                      \includegraphics[width=0.9\linewidth]{logo}

                  \end{itemize}
                \end{column}
              \end{columns}
            \end{block}
            \vfill
            \begin{block}{Results 3}
              \begin{itemize}
              \item Another block with more info
              \item One of this blocks should produce the ``Oh!'' effect
              \item Graphs, charts, plots will improve the quality
              \end{itemize}
            \end{block}
            \vfill
            \begin{block}{Conclusions}
              \begin{itemize}
              \item Point out what are the most important conclusions of your work
              \item What is new in your work
              \item What is better than older methods
              \item And maybe what could still be improved
              \item As you see, the poster main parts are:
                \begin{itemize}
                \item Abstract
                \item Introduction presenting context and problem
                \item Solution proposed
                \item Results
                \item Conclusions
                \end{itemize}
              \item Remember: a good poster prefers visual information over text
              \end{itemize}
            \end{block}
          }
          % ---------------------------------------------------------%
          % end the column
        \end{minipage}
      \end{beamercolorbox}
    \end{column}
    % ---------------------------------------------------------%
    % end the column
  \end{columns}
  \vskip1ex
  %\tiny\hfill\textcolor{ta2gray}{Created with \LaTeX \texttt{beamerposter}  \url{http://www-i6.informatik.rwth-aachen.de/~dreuw/latexbeamerposter.php}}
  \tiny\hfill{Created with \LaTeX~\texttt{beamerposter}  \url{http://www-i6.informatik.rwth-aachen.de/~dreuw/latexbeamerposter.php} \hskip1em}
\end{frame}
\end{document}


%%%%%%%%%%%%%%%%%%%%%%%%%%%%%%%%%%%%%%%%%%%%%%%%%%%%%%%%%%%%%%%%%%%%%%%%%%%%%%%%%%%%%%%%%%%%%%%%%%%%
%%% Local Variables: 
%%% mode: latex
%%% TeX-PDF-mode: t
%%% End:
